\documentclass[a4paper]{article}

\usepackage{amsmath,amsthm,amssymb, titling, tikz, mdframed}

\title{\vspace{-2cm}Lecture X\vspace{-2cm}}
\date{}

\newmdtheoremenv{theorem}{Theorem}
\newmdtheoremenv{definition}{Definition}
\newmdtheoremenv{corollary}{Corollary}
\newmdtheoremenv{algorithm}{Algorithm}
\newmdtheoremenv{lemma}{Lemma}

\begin{document}
\maketitle
\section{Exercises}
\noindent \textbf{Exercise 10.1}  Analyze the correctness of the WITNESS protocol.

\section{Lecture Summary}
\begin{algorithm}
\textbf{Randomized Communication Protocol WITNESS -} \\
Let $\operatorname{PRIM}(N)$ denote the set of all prime numbers not exceeding $N$, and let $\operatorname{num}(x)$ denote the integer whose binary representation is the bit string $x$. \\
\begin{tabular}{c c}
Given: & Computer $C_{1}$ with $n$ bits $x_{1}x_{2}\cdots x_{n}$ \\
& Computer $C_{2}$ with $n$ bits $y_{1}y_{2}\cdots y_{n}$
\end{tabular} \\
\\
\textbf{Phase 1:} $C_{1}$ chooses uniformly at random a prime $p$ from $\operatorname{PRIM}(n^{2})$. \\
\\
\textbf{Phase 2:} $C_{1}$ computes the residue
\begin{center}
$s \equiv \operatorname{num}(x) \pmod{p}$
\end{center}
and sends $s$ and $p$ in binary representation to $C_{2}$.
\\
\textbf{Phase 3:} After receiving $s$ and $p$, computer $C_{2}$ computes
\begin{center}
$q \equiv \operatorname{num}(y) \pmod{p}.$
\end{center}
If $q \neq s$, then $C_{2}$ outputs "$x$ and $y$ are not equal." \\
If $q = s$, then $C_{2}$ outputs "$x$ and $y$ are equal."
\end{algorithm}
\end{document}
